% Independent exact derivation. Prefix HSQ. Requires amsmath,amssymb,hyperref.
\section{Every actual off-real heat quartet and its full supported jet receiver}
\label{sec:full-support-heat-quartet}

\paragraph{Sources and scope.}
The unchanged kernel and heat coordinate are Brad Rodgers and Terence
Tao's \href{https://arxiv.org/abs/1801.05914v5}{\emph{The De
Bruijn--Newman constant is non-negative}}, original author source
\nolinkurl{fmp-template.tex}, equations \texttt{hoz}, \texttt{sas},
\texttt{phidef}, and \texttt{htdef}, lines 78--103. Source lines
74--110 were read for this calculation. The original source's
SHA--256 is
\begin{center}\small
\nolinkurl{25b015c0fb151f1613247d82cf53a6561a320d4dab2bb4f809482e7c477a5817}.
\end{center}
The complete included theta heat proof THS1--31 and the complete
supported divisor proof HZD1--18 were read.
Their integral bounds, actual divisor, and full coefficient map are
retained below with their exact constants. The split-support carrier
is the existing construction of The Clankers,
\href{https://github.com/KokunoYumeto/zeta-function-research-reader/blob/a2851f9eb352e7e4376bc629de86fa4ed9c1c434/workbenches/splitzero-tandem/foundations/split-support-geometry-v11/split_support_geometry_arithmetic_curve_v11.tex#L1226}{\emph{Split
Support Geometry}, v11, Definition \texttt{def:lattice-split}}.
All finite-jet quotient, symmetry, support, and exclusion statements
made here are proved below. The construction is indexed by actual
zeros of the given function; it does not assume that the indexing
set is nonempty.

\subsection{The original moments and the exclusion inequality}

For $u\ge0$ and real $t$ retain exactly
\begin{equation}\tag{HSQ1}
 \begin{gathered}
 \Phi(u)=\sum_{n\ge1}
       (2\pi^2n^4e^{9u}-3\pi n^2e^{5u})e^{-\pi n^2e^{4u}},\\
 H_t(z)=\int_0^\infty e^{tu^2}\Phi(u)\cos(zu)\,du.
 \end{gathered}
\end{equation}
Every summand of the kernel is positive: its prefactor is
$\pi n^2e^{5u}(2\pi n^2e^{4u}-3)>0$. THS3--6 prove convergence,
all differentiations used below, and the explicit estimates
\begin{equation}\tag{HSQ2}
 \begin{gathered}
 0<\pi(2\pi-3)e^{9u-\pi e^{4u}}
       \le\Phi(u)\le C_\Phi e^{9u-\pi e^{4u}},\\
 a=e^{-3\pi},\qquad
 C_\Phi=2\pi^2\frac{1+11a+11a^2+a^3}{(1-a)^5},\\
 \mathcal M_p(A,B)=\frac{C_\Phi}{2\pi}
  \exp\left(\frac{3A^2}{32\pi}
      +\frac{(B+9+p)^2}{8\pi}-\frac\pi2\right),\qquad A,B\ge0.
 \end{gathered}
\end{equation}
In particular the integral of $u^pe^{Au^2+Bu}\Phi(u)$ is at
most $\mathcal M_p(A,B)$. For a direct convergence check of this
bound, use $u^p\le e^{pu}$, reserve half of
$-\pi e^{4u}$, and apply $e^{4u}\ge(32/3)u^4$ and
$e^{4u}\ge8u^2$ to the other two quarters. Completing squares
bounds the remaining positive exponent by
$3A^2/(32\pi)+(B+9+p)^2/(8\pi)$. Finally
$e^{4u}\ge1+4u$ bounds the integral of the reserved exponential
by $e^{-\pi/2}/(2\pi)$. This recovers the exact constant in
(HSQ2) and justifies every moment integral below.

For every real $y$ and integer $r\ge0$, define the actual positive
moments
\begin{equation}\tag{HSQ3}
 \begin{gathered}
 M_{2r}(t,y)=\int_0^\infty u^{2r}e^{tu^2}\Phi(u)\cosh(yu)\,du
              =\partial_t^rH_t(iy),\\
 c_t=\pi(2\pi-3)\exp\bigl(\min(t,0)-\pi e^4\bigr),\qquad
 \frac{c_t}{2r+1}\le M_{2r}(t,y)
        \le\mathcal M_{2r}(\max(t,0),|y|).
 \end{gathered}
\end{equation}
The upper bound uses $\cosh(yu)\le e^{|y|u}$ and (HSQ2).
For the lower bound restrict to $0\le u\le1$, use
$\cosh(yu)\ge1$, $e^{tu^2}\ge e^{\min(t,0)}$, and the lower
kernel bound in (HSQ2), then integrate $u^{2r}$. Thus
$H_t(iy)=M_0(t,y)>0$ and $\partial_tH_t(iy)=M_2(t,y)>0$
for all the stated real coordinates.
In terms of the original coefficients $h_j(t)$ of THS8 the same
moments are exactly
\begin{equation}\tag{HSQ4}
 M_{2r}(t,y)=\sum_{j=0}^{\infty}
   \frac{(-1)^{j+r}(2j+2r)!}{(2j)!}\,h_{j+r}(t)y^{2j}.
\end{equation}
Indeed expand $\cosh(yu)$ in its power series and substitute the
moment definition of $h_{j+r}$. Absolute convergence is bounded
by (HSQ3); all its moment summands are nonnegative for real $t$.
This proves the formula with every original factorial and sign.

For every real $x,y,t$, the full real and imaginary parts are
\begin{equation}\tag{HSQ5}
 \begin{split}
 \operatorname{Re}H_t(x+iy)
   &=\int_0^\infty e^{tu^2}\Phi(u)\cos(xu)\cosh(yu)\,du,\\
 \operatorname{Im}H_t(x+iy)
   &=-\int_0^\infty e^{tu^2}\Phi(u)\sin(xu)\sinh(yu)\,du.
 \end{split}
\end{equation}
The elementary inequality $\cos v\ge1-v^2/2$ is strict for
$v\ne0$. To verify strictness, the even function
$\cos v-1+v^2/2$ vanishes at zero and has derivative
$v-\sin v>0$ for $v>0$: the latter is the integral of
$1-\cos s$, positive on any such interval. Integrating the
inequality against the positive weight in (HSQ5) gives
\begin{equation}\tag{HSQ6}
 \operatorname{Re}H_t(x+iy)
   \ge H_t(iy)-\frac{x^2}{2}\partial_tH_t(iy).
\end{equation}
For $x\ne0$ this is a strict inequality, because its exact
remainder is the integral of the strictly positive function
$e^{tu^2}\Phi(u)\cosh(yu)
 [\cos(xu)-1+x^2u^2/2]$ on $u>0$. Its integrability follows
from (HSQ3) for $r=0,1$.
Consequently every zero of $H_t$ at $x+iy$ satisfies
\begin{equation}\tag{HSQ7}
 x\ne0,\qquad
 x^2\partial_tH_t(iy)>2H_t(iy)>0,
 \qquad
 x^2>\frac{2M_0(t,y)}{M_2(t,y)}
       \ge\frac{2c_t}{\mathcal M_2(\max(t,0),|y|)}.
\end{equation}
In particular the requested non-strict necessary bound follows,
and the closed region $x^2\le2M_0(t,y)/M_2(t,y)$ is actually
zero free. A completely explicit uniform version is obtained for
$t_-\le t\le t_+$ and $|y|\le Y$ by replacing its last ratio
by $2c_{t_-}/\mathcal M_2(\max(t_+,0),Y)$. No zero location,
time, or kernel has been changed in deriving these bounds.

\subsection{The four distinct original points and their common multiplicity}

Let $Z_t^{\mathrm{off}}=\{z\in\mathbb C\setminus\mathbb R:
H_t(z)=0\}$ be the actual zero set at a fixed real $t$.
For each $z_0=x+iy\in Z_t^{\mathrm{off}}$, both $x$ and $y$
are nonzero: $y\ne0$ by its definition, and $x\ne0$ by
(HSQ3). Retain the ordered four points
\begin{equation}\tag{HSQ8}
 w_1=z_0=x+iy,\quad w_2=-z_0=-x-iy,\quad
 w_3=\overline{z_0}=x-iy,\quad
 w_4=-\overline{z_0}=-x+iy.
\end{equation}
They are distinct. For example an equality with its negative
would force $x=y=0$, an equality with its conjugate would force
$y=0$, and an equality with its negative conjugate would force
$x=0$; these cover the pairwise possibilities after symmetry.
They are the four corners of the actual rectangle with horizontal
coordinates $\pm x$ and vertical coordinates $\pm y$.

The source integral gives $H_t(-z)=H_t(z)$ and
$H_t(\bar z)=\overline{H_t(z)}$. Its derivatives obey
\begin{equation}\tag{HSQ9}
 H_t^{(k)}(-z)=(-1)^kH_t^{(k)}(z),\qquad
 H_t^{(k)}(\bar z)=\overline{H_t^{(k)}(z)}.
\end{equation}
The first follows by differentiating the even identity, the second
by conjugating its local power series. Since $H_t(0)>0$, this
entire function is not identically zero. Its first nonzero Taylor
coefficient at any root therefore has a finite positive index.
Equations (HSQ9) prove that the exact multiplicity is one and the
same integer $m\ge1$ at all four points. In particular the four
leading coefficients remain
$H_t^{(m)}(w_i)/m!$ with the signs and conjugations in (HSQ9).
None is replaced by one.

\subsection{The exact polynomial and the entire residual factor}

Define the monic polynomial with these actual coordinates
\begin{equation}\tag{HSQ10}
 \begin{split}
 p_Q(z)&=\prod_{i=1}^4(z-w_i)^m\\
 &=\left[(z^2-z_0^2)(z^2-\bar z_0^{\,2})\right]^m\\
 &=\left[z^4-2(x^2-y^2)z^2+(x^2+y^2)^2\right]^m,
 \qquad P_i(z)=\prod_{j\ne i}(z-w_j)^m.
 \end{split}
\end{equation}
All three forms follow by multiplying the specified factors; the
first form retains the four separately indexed points. Division
gives an actual entire residual factor
\begin{equation}\tag{HSQ11}
 H_t(z)=p_Q(z)U_Q(z),\qquad
 U_Q(w_i)=\frac{H_t^{(m)}(w_i)}{m!P_i(w_i)}\ne0.
\end{equation}
Away from the four points take the ordinary quotient. At each
point its Taylor factor $(z-w_i)^m$ cancels and the other factors
are units, extending it holomorphically. This proves (HSQ11)
globally. In the original coordinates its denominators are
\begin{equation}\tag{HSQ12}
 (P_1(w_1),P_2(w_2),P_3(w_3),P_4(w_4))
 =\bigl((8ixyz_0)^m,(-8ixyz_0)^m,
        (-8ixy\bar z_0)^m,(8ixy\bar z_0)^m\bigr).
\end{equation}
They follow by multiplying the three original differences at
each point. All are nonzero by (HSQ8). The residual factor can
have other zeros; no global inverse for it is asserted.

\subsection{Explicit Hermite idempotents and the complete CRT map}

Put $\Omega=\mathbb C\setminus\mathbb R$ and
$A=\mathcal O(\Omega)$, keeping both components. For every $i$
define, exactly in the original variable,
\begin{equation}\tag{HSQ13}
 Q_i(z)=\sum_{k=0}^{m-1}
      \frac{(1/P_i)^{(k)}(w_i)}{k!}(z-w_i)^k,
 \qquad E_i(z)=P_i(z)Q_i(z).
\end{equation}
The coefficient of $(z-w_i)^k$ can equivalently be written as
\begin{equation}\tag{HSQ14}
 \frac1{P_i(w_i)}
 \sum_{\substack{\alpha_j\ge0\ (j\ne i)\\\sum_{j\ne i}\alpha_j=k}}
 \prod_{j\ne i}
   \frac{(-1)^{\alpha_j}\binom{m+\alpha_j-1}{\alpha_j}}
        {(w_i-w_j)^{\alpha_j}}.
\end{equation}
Indeed expand each factor
$(1+(z-w_i)/(w_i-w_j))^{-m}$ near $w_i$, using its convergent
negative-integer binomial series, and multiply the three series.
Only finitely many terms contribute to any displayed coefficient.
This proves (HSQ14) with its original differences and signs.

By Taylor multiplication, $E_i\equiv1\pmod{(z-w_i)^m}$;
at every other $w_j$ its factor $P_i$ gives
$E_i\equiv0\pmod{(z-w_j)^m}$. Each $E_i$ has degree less than
$4m$. A polynomial vanishing to order at least $m$ at all four
distinct points is divisible by their product $p_Q$: successive
division by the individual factors preserves the remaining
vanishing orders since their values at other points are nonzero.
It follows that
\begin{equation}\tag{HSQ15}
 \sum_iE_i=1\quad\hbox{as polynomials},\qquad
 E_iE_j\equiv\delta_{ij}E_i\pmod{p_Q}.
\end{equation}
For the first identity the difference has degree less than $4m$
and is divisible by the degree-$4m$ monic polynomial. The second
follows by checking all four jets and the same divisibility test.

Retain four distinct formal coordinates and rings
$D_i=\mathbb C[\epsilon_i]/(\epsilon_i^m)$ and $B_Q=\prod_iD_i$.
The explicit ring map and its inverse on the quotient are
\begin{equation}\tag{HSQ16}
 \begin{gathered}
 J_Q:A/(p_Q)\xrightarrow{\ \sim\ }B_Q,\qquad
 [f]\longmapsto
  \left(\sum_{k=0}^{m-1}\frac{f^{(k)}(w_i)}{k!}\epsilon_i^k\right)_{i=1}^4,\\
 (b_i(\epsilon_i))_i\longmapsto
        \left[\sum_{i=1}^4E_i(z)b_i(z-w_i)\right],\qquad
 J_Q([z])=(w_i+\epsilon_i)_i.
 \end{gathered}
\end{equation}
The product rule proves multiplicativity of each jet map with
its stated factorials. The displayed inverse realizes any tuple
by a polynomial, because of (HSQ15). The kernel on $A$ is exactly
$(p_Q)$: a vanishing jet makes $f/p_Q$ extend holomorphically
at its corresponding point, and away from the four points it is
already holomorphic on both components of $\Omega$. Conversely
a multiple of $p_Q$ has all four zero jets. This proves all maps
and every claimed equality in (HSQ16).

Now retain the original heat divisor ring $R_t=A/(H_t|_\Omega)$.
Equation (HSQ11) proves $(H_t)\subseteq(p_Q)$, so its actual
quartet receiver is the surjection
\begin{equation}\tag{HSQ17}
 \pi_Q:R_t\twoheadrightarrow B_Q,\qquad
 [f]\longmapsto\left(\sum_{k=0}^{m-1}
       \frac{f^{(k)}(w_i)}{k!}\epsilon_i^k\right)_i,
 \qquad \ker\pi_Q=(p_Q)/(H_t).
\end{equation}
Surjectivity and the kernel follow from (HSQ16). The Hermite
polynomial is a section of $A\to A/(p_Q)$ as a finite-dimensional
complex vector space, with subsequent remainder understood in
the quotient. It is not being asserted to define a multiplicative
section $B_Q\to R_t$. The original $H_t$, the actual residual
factor $U_Q$, and the exact kernel in (HSQ17) remain present.

\subsection{The whole split-support receiver and the complete coefficient map}

Let $L$ be the full nontrivial bounded distributive lattice and
write $S_T=G_L(T)$ for any ring $T$. Retain
\begin{equation}\tag{HSQ18}
 \begin{gathered}
 \mathscr D_{t,L}=G_L(R_t),\qquad
 \mathscr Q_{t,z_0,L}=G_L(B_Q),\qquad
 G_L(\pi_Q)([f],\lambda)=(\pi_Q[f],\lambda),\\
 \mathscr D_{t,L}/\langle([p_Q],1_L)\sim(0,1_L)\rangle
                \simeq\mathscr Q_{t,z_0,L}.
 \end{gathered}
\end{equation}
The full operations are
$(b,\lambda)+(c,\mu)=(b+c,\lambda\vee\mu)$ and
$(b,\lambda)(c,\mu)=(bc,\lambda\wedge\mu)$, with nonzero
amplitudes allowed only at top support. The map is a surjective
semiring homomorphism by (HSQ17), with every label unchanged.
Its kernel congruence is equality of labels and amplitude
difference in $\ker\pi_Q$. At top support every difference
is a multiple of $[p_Q]$, so multiplication of the displayed
generating relation followed by addition generates that entire
congruence. Below top support the amplitude is already zero.
This proves the supported quotient, including its zero convention.

Let $K$ be the retained tropical semifield with zero, let
$K_L=K\otimes_{\mathbb B}L$, and keep
$\iota_L(\lambda)=1_K\otimes\lambda$ with its exact retraction
$\rho_L(\sum a_i\otimes\lambda_i)=\bigvee_{a_i\ne0}\lambda_i$.
The complete maps are
\begin{equation}\tag{HSQ19}
 \Phi_T:S_T\hookrightarrow T\times K_L,
 \qquad(b,\lambda)\longmapsto(b,\iota_L(\lambda)),\qquad
 \Phi_{B_Q}G_L(\pi_Q)
       =(\pi_Q\times\operatorname{id}_{K_L})\Phi_{R_t}.
\end{equation}
Both coordinates preserve the operations, and the retraction
of $\iota_L$ proves injectivity. The equality is verified on
each displayed pair. Thus the complete map carries the actual
quartet quotient, with all four jets and its original shared label.
In particular the image of the supported equation $H_t=0$ is
$(0_{B_Q},\iota_L(1_L))$, not $(0_{B_Q},0_{K_L})$.

For each $i$ let $q_i\in B_Q$ be the element with unit in its
$i$th ring component and zero in the other three; it is the class
of $E_i$. Its full-supported lift satisfies
\begin{equation}\tag{HSQ20}
 \widehat q_i^{\,2}=\widehat q_i,\qquad
 \widehat q_i\widehat q_j=e\quad(i\ne j),\qquad
 \sum_i\widehat q_i=1_{G_L(B_Q)}.
\end{equation}
The middle amplitude is zero but its meet of top labels is top.
Each coordinate $\epsilon_i$ retains order $m$ in its ring factor;
the lift of its zero $m$th power is likewise the supported zero.
The receiver has complex ring dimension $4m$, with the original
basis $q_i\epsilon_i^k$, $1\le i\le4$, $0\le k<m$.

There are exactly four arithmetic primes in this receiver, namely
\begin{equation}\tag{HSQ21}
 \mathfrak Q_i=\{(0,\lambda):\lambda\in L\}
       \cup\{(b,1_L):b_i(0)=0\}\qquad(1\le i\le4).
\end{equation}
To prove exhaustion on the ring side, the orthogonal idempotents
$q_i$ sum to one, so a prime selects exactly one factor: two
excluded idempotents would have zero product, and excluding none
would put their sum in the prime. In that factor $\epsilon_i$
is nilpotent and hence belongs to every prime; quotienting by it
gives $\mathbb C$, which has only the zero prime. The full
split-support classification therefore adds exactly the support
primes $\{(0,\lambda):\lambda\in P\}$ for
$P\in\operatorname{SpecLat}L$, each contained in all four
(HSQ21). The four primes are distinct, as detected by the
specified $q_i$. None has been identified with another by a symmetry.

\subsection{Shared labels, independent labels, and the exact synchronization}

The canonical comparison is the injective unital semiring map
\begin{equation}\tag{HSQ22}
 \theta:G_L\!\left(\prod_iD_i\right)\hookrightarrow
           \prod_iG_L(D_i),\qquad
 ((b_i)_i,\lambda)\longmapsto((b_i,\lambda))_i.
\end{equation}
Its image consists exactly of the tuples whose four labels are
equal. The individual carrier conditions then imply that any
nonzero amplitude makes that common label top. Conversely those
conditions are exactly the original carrier of $G_L(B_Q)$.
All operations and the injection follow coordinatewise.

The strongest elementary synchronization is the actual retraction
\begin{equation}\tag{HSQ23}
 r:\prod_iG_L(D_i)\longrightarrow G_L(B_Q),\qquad
 ((b_i,\lambda_i))_i\longmapsto
       ((b_i)_i,\bigvee_i\lambda_i),\qquad r\theta=\operatorname{id}.
\end{equation}
A nonzero ring tuple forces at least one $\lambda_i$ to be top,
so this belongs to the required carrier. The retraction preserves
addition, zero, and unit. More precisely it is a homomorphism of
$G_L(B_Q)$-semimodules when the product is acted on through
$\theta$: for a scalar label $\nu$, distributivity gives
$\bigvee_i(\nu\wedge\lambda_i)=\nu\wedge\bigvee_i\lambda_i$,
and the amplitude equality is componentwise multiplication.
Thus $\theta r$ is an idempotent linear projection onto the
shared-label object, with every amplitude coordinate untouched.

It is not multiplicative in general. The entire defect is explicit:
for $v=((b_i,\lambda_i))$, $w=((c_i,\mu_i))$, put
\begin{equation}\tag{HSQ24}
 \alpha=\bigvee_i(\lambda_i\wedge\mu_i),\qquad
 \beta=(\bigvee_i\lambda_i)\wedge(\bigvee_j\mu_j).
\end{equation}
Then $\alpha\le\beta$ and
$r(vw)=((b_ic_i)_i,\alpha)$ while
$r(v)r(w)=((b_ic_i)_i,\beta)=r(vw)+z_\beta$.
Taking zero amplitudes, top label only in the first coordinate
of $v$ and only in the second of $w$, gives $\alpha=0_L$ and
$\beta=1_L$, so the defect is real.

In fact no unital semiring retraction of (HSQ22) exists for this
nontrivial $L$. For a direct proof let $v_i$ be the product element
with unit in coordinate $i$ and global zero in the other coordinates.
For $i\ne j$, $v_iv_j=0$, whereas
$v_i+\theta(e)=\theta(\widehat q_i)$. If a semiring retraction
$s$ existed, the latter identity would give
$s(v_i)+e=\widehat q_i$. Its amplitude is the nonzero ring
element $q_i$, forcing $s(v_i)=\widehat q_i$ by the carrier
condition. It follows that $s(v_iv_j)=\widehat q_i\widehat q_j=e$,
contradicting preservation of zero, since $e\ne\tau$. The
module retraction (HSQ23) is therefore retained in its proved
category, rather than promoted to a nonexistent semiring retraction.

The full coefficient maps participate in exactly the same relation.
Define
\begin{equation}\tag{HSQ25}
 \begin{gathered}
 \Gamma:B_Q\times K_L\hookrightarrow\prod_i(D_i\times K_L),
 \qquad((b_i),k)\longmapsto((b_i,k))_i,\\
 \mathcal R:\prod_i(D_i\times K_L)\longrightarrow B_Q\times K_L,
 \qquad((b_i,k_i))_i\longmapsto((b_i),\sum_i k_i),\\
 (\prod_i\Phi_{D_i})\theta=\Gamma\Phi_{B_Q},\qquad
 \mathcal R(\prod_i\Phi_{D_i})=\Phi_{B_Q}r,
 \qquad\mathcal R\Gamma=\operatorname{id}.
 \end{gathered}
\end{equation}
The map $\Gamma$ is a semiring embedding with common second
coordinate as its image. The map $\mathcal R$ is additive and
linear for the $B_Q\times K_L$-action through $\Gamma$.
Distributivity and the idempotence of addition in $K_L$ prove its
linearity and retraction property. For the two remaining equations
use $\sum_i\iota_L(\lambda_i)=\iota_L(\bigvee_i\lambda_i)$.
On the original supported image, the common tropical value is
exactly $\iota_L(\lambda)$, retaining all of $L$. These are the
complete maps, not a projection to a chosen lattice character.

\subsection{Exact symmetry actions on the four jet factors}

On $A$ define $\mathsf a f(z)=f(-z)$ and
$\mathsf c f(z)=\overline{f(\bar z)}$. The first is a complex
linear ring involution; the second is a conjugate-linear ring
involution. They commute, preserve $H_t$ and $p_Q$, and hence act
on both quotients in (HSQ17). Write
$a=(1\ 2)(3\ 4)$ and $c=(1\ 3)(2\ 4)$ for their point
permutations. In the exact jet coordinates, their actions are
\begin{equation}\tag{HSQ26}
 \begin{split}
 (\mathsf a b)_i(\epsilon_i)
   &=\sum_{k=0}^{m-1}(-1)^k b_{a(i),k}\epsilon_i^k,\\
 (\mathsf c b)_i(\epsilon_i)
   &=\sum_{k=0}^{m-1}\overline{b_{c(i),k}}\epsilon_i^k,
 \qquad b_i(\epsilon_i)=\sum_{k=0}^{m-1}b_{i,k}\epsilon_i^k.
 \end{split}
\end{equation}
Differentiate $f(-z)$ $k$ times and evaluate at $w_i$.
The result is $(-1)^kf^{(k)}(-w_i)$, proving the first formula.
For the second, conjugate the Taylor expansion at $\bar w_i$.
These also prove compatibility with (HSQ16)--(HSQ17), including
every coefficient sign and conjugation. Squaring either formula
gives the identity; composing them gives the remaining involution
with coefficients $(-1)^k\overline{b_{ac(i),k}}$, where
$ac=(1\ 4)(2\ 3)$.

For the original coordinate element $Z=(w_i+\epsilon_i)_i$,
$\mathsf aZ=-Z$ and $\mathsf cZ=Z$. The second equality follows
because this conjugate-linear operation includes conjugation of
the coefficient and permutes its centre; it is not coefficientwise
conjugation at a fixed centre. The Hermite polynomials themselves
satisfy the exact polynomial identities
\begin{equation}\tag{HSQ27}
 E_i(-z)=E_{a(i)}(z),\qquad
 \overline{E_i(\bar z)}=E_{c(i)}(z).
\end{equation}
Both sides have degree less than $4m$ and exactly the same four
jets, by (HSQ13)--(HSQ15); their difference is divisible by $p_Q$
and is therefore zero. Thus these symmetry assertions preserve the
explicit CRT inverse as well as the quotient map.

Every symmetry lifts to $G_L$ by changing its amplitude as in
(HSQ26) and leaving its single support label unchanged. It
commutes with $\Phi$ when the same amplitude action and the
identity on $K_L$ are used on the complete receiver. On the
independently labelled product it also permutes the four labels
with its factor permutation. Their join is unchanged, so both
synchronizations in (HSQ23) and (HSQ25) commute with the actions.
At no stage is the invariant subring substituted for $B_Q$, nor
are orbit points or their nilpotent coordinates identified.
The receiver therefore retains four distinct arithmetic points,
all $4m$ jet coordinates, every lattice label, and the original
coordinate constraints (HSQ7)--(HSQ8).
